\chapter{Electrocardiogram}

\section*{Overview}
An electrocardiogram (ECG or EKG) records the heart's electrical activity through electrodes placed on the skin.

\section*{Why It Is Done}
It evaluates chest pain, palpitations, fainting, shortness of breath, arrhythmias, heart attack, conduction problems, and effects of medicines or electrolyte abnormalities.

\section*{How It Is Performed}
The patient lies still while adhesive electrodes are placed on the chest and limbs. A resting tracing usually takes only a few minutes; ambulatory and exercise ECGs use related monitoring methods.

\section*{Risks and Results}
A resting ECG is noninvasive and has no significant medical risk, although removing electrodes may irritate the skin. A clinician interprets rhythm, rate, intervals, conduction, and signs of ischemia or prior injury.

\section*{References}
\begin{enumerate}
  \item MedlinePlus. Electrocardiogram. U.S. National Library of Medicine; 2025.
  \item Current Medical Diagnosis and Treatment. McGraw-Hill; 2025.
\end{enumerate}
\clearpage
